\section{Consequences for the Linnik constant and for RH}\label{gl:sec:K}

\subsection{The Linnik constant}

\begin{corollary}\label{gl:cor:K7}
Suppose that for every $q\ge3$ the function $\varphi_{\mathbb Q(\zeta_q),0}$ is the Laplace transform of a GGC. Then every sufficiently large even integer is a sum of two primes and at most seven powers of two.
\end{corollary}

% CHECK: Theorem 2 of HBP02 (K=7 under GRH) is confirmed against arXiv math/0201299v2 in misc/reviews/nick-linnik/hbp02-check.md, and the proof cites it without the source's claims that step (3) is "the only GRH input" and that their theorem "applies verbatim". Lemma 7 of that paper, the major arc lower bound with the constant 2.7895, is stated for K at least 9, and the paper does not say how it is used for K=7 (question for the authors, misc/catalog/QUESTIONS.md). K=7 under GRH for PR03 is of medium confidence there.
\begin{proof}
Theorem~\ref{gl:thm:grh} gives GRH for all Dirichlet $L$ functions, and under GRH every sufficiently large even integer is a sum of two primes and at most seven powers of two, by the theorem of Heath-Brown and Puchta~\cite{HBP02} (see also Pintz and Ruzsa~\cite{PR03}).
\end{proof}

The corollary is weaker than the known conditional result. By Corollary~\ref{gl:cor:basepoints}(a) and equation~\eqref{gl:eq:cyclotomic}, its hypothesis is GRH for every Dirichlet $L$ function together with $L(\tfrac12,\chi)\ne0$ for every primitive character $\chi$. It contains GRH and adds a statement that is not known to follow from it, because the nonvanishing at the centre is open (for real characters it is Chowla's conjecture), whereas $K=7$ already follows from GRH alone. A direct proof of the Thorin property for all cyclotomic fields would prove GRH, and with it $K=7$ unconditionally, one better than the record $K=8$ of Pintz and Ruzsa~\cite{PR20}.

For $K$ itself the Thorin property gives what GRH gives. Nothing in the proof of Corollary~\ref{gl:cor:K7} uses the nonvanishing at the centre, and the inputs that decide $K$ below seven, the large deviation estimate for the sum $\mathcal D(\alpha)$ over the powers of two (step~(4) of Section~\ref{gl:sec:where}), the sieve constant for prime differences (step~(2)) and the refinements of Pintz and Ruzsa, are additive questions about the powers of two and not statements about zeros. With the present methods the framework therefore does not reach below seven.

A weaker hypothesis may suffice for seven. The Thorin property at tilt $u_0=1$ for cyclotomic fields (Corollary~\ref{gl:cor:basepoints}(b) with equation~\eqref{gl:eq:cyclotomic}) is GRH with possible real zeros. Where could a real zero hurt? Write $\mathcal S(a/q+\eta)$ as a combination of twisted prime sums over the characters modulo $q$. A real zero $\beta$ of $L(s,\chi^*)$, with $\chi^*$ of conductor $q^*\mid q$, contributes a term of size about $q^{1/2}\varphi(q)^{-1}N^\beta(1+|\eta|N)^{-1}$, which exceeds $N^{3/4}$ only if $q<N^{2\beta-3/2}$ and $|\eta|<N^{\beta-7/4}$. Those arcs lie close to the major arcs, and moving them there, where a real zero changes the main term by a factor of the shape $1-\chi^*(\cdot)N^{\beta-1}$ that Gallagher's method~\cite{Gallagher75} already handles for one exceptional zero, would reduce the question to major arc estimates with real zeros (Problem~\ref{gl:op:exceptional}). If it works, real zeros, and with them Siegel zeros, become irrelevant to $K=7$.

\subsection{What would help towards RH}\label{gl:sec:rh}

The three smoothings of the centre clock are compared in Chapter~\ref{chap:tf}.

A proof has to establish a closed positivity condition at the boundary, from input that sees the arithmetic beyond positivity. Zero free regions of the kind announced in October 2026~\cite{OAI26a} would be such input, if confirmed, and each one would move the unconditional end of the basepoint family inwards (Section~\ref{gl:sec:oct}); averages over families of $L$ functions, where the low lying zeros are conjectured to obey the symmetry laws of Katz and Sarnak~\cite{KatzSarnak99}, are another (Problem~\ref{gl:op:negmoment}). Section~\ref{tf:sec:past} says what would get past each of the three obstructions of Chapter~\ref{chap:tf}.
